\section{Related Work}
\label{sec:related}

\paragraph{Diffusion-based image-to-video generation.}
With the recent advances in diffusion models~\citep{DiffusionModels,DDPM}, image animation has achieved tremendous progress~\citep{wang2023videocomposer,VideoCrafter1,VideoCrafter2,AnimateDiff}.
Early methods inflate pre-trained text-to-image models~\citep{LatentDiffusion} to add motions to an image~\citep{Tune-a-video,Text2VideoZero,Make-A-Video}.
A notable example is AnimateDiff~\citep{AnimateDiff}, which learns low-rank adapters~\citep{LoRA}  for different motions.
Later works seek to inject a conditioning frame into a pre-trained text-to-video model~\citep{VideoDiffusionModels,ImagenVideo,LVDM}.
VideoCrafter1~\citep{VideoCrafter1} leverages a dual cross-attention layer to condition on features of both the image and the text prompt.
DynamicCrafter~\citep{DynamiCrafter} further improves it by concatenating the input image with noisy latent.
Stable Video Diffusion (SVD)~\citep{blattmann2023stable} instead works in an image-only manner, which replaces the CLIP~\citep{CLIP} text embedding of the text prompt with the CLIP image embedding of the conditioning frame.
However, none of these models support direct trajectory control of scene elements, and rather require multiple attempts to obtain a desired result.
In this work, we aim to enable intuitive motion control in animating a pre-existing image.
Since SVD only takes in an image without any text prompt, we utilize it as the base model following prior work~\citep{yin2023dragnuwa,wu2024draganything}.





\paragraph{Spatial control in image diffusion models.}
One common way to incorporate spatial control into the image generation process is to fine-tune pre-trained models to incorporate conditioning on depth maps or bounding boxes~\citep{ControlNet,ye2023ip,SpaText,GLIGEN,PairDiffusion,InstanceDiffusion}.
While these methods demonstrate high fidelity, they require excessive computing resources and labor-intensive data annotations. 
Therefore, several tuning-free approaches have been proposed~\citep{MasaCtrl,LayoutGuidanceAttn1,LayoutGuidanceAttn2,PromptToPrompt}.
Self-Guidance~\citep{epstein2023diffusion} Attend-and-Excite~\citep{chefer2023attend} , and TraDiffusion~\citep{wu2024tradiffusion} control the image layout by manipulating intermediate attention maps. %
They first estimate the generated objects' positions using the attention maps produced from text--image cross-attention layers and then optimize the latent to increase the attention values at specific positions.
However, these approaches require associating the control target with a specific token in the text prompt and thus cannot be directly applied to image-only editing.
Closest to ours are methods that alter image layouts without text input~\citep{pan2023drag,mou2023dragondiffusion,mou2024diffeditor,ling2023freedrag,liu2024drag,zhang2024gooddrag,cui2024stabledrag,hou2024easydrag,zhao2024fastdrag,shi2024dragdiffusion}.
They enforce similarity between semantically correlated feature maps extracted from the upsampling blocks of a denoising U-Net~\citep{tang2023emergent,namekata2024emerdiff,hedlin2024unsupervised,zhang2023tale,luo2023dhf}.
In contrast, we show that the cross-frame semantic correspondence of upsampling feature maps in image-to-video diffusion models~\citep{blattmann2023stable} is weak. Thus, optimization directly based on these feature maps leads to sub-optimal results.


\paragraph{Motion control in video diffusion models.}
Several recent works have studied camera pose control in video diffusion models~\citep{CameraCtrl,xu2024camco,kuang2024collaborative,MotionMaster,xiao2024video,hou2024training,bahmani2024vd3d,li2024image,zhang2024dragentity}.
The representative work MotionCtrl~\citep{MotionCtrl} fine-tunes pre-trained video generators to follow camera trajectory input. 
However, video datasets with accurate camera pose annotations are limited~\citep{RealEstate10k}, and fine-tuning video models requires high computation costs.
Another line of work focuses on object trajectory control~\citep{wang2023videocomposer,yin2023dragnuwa,wu2024draganything,zhou2024trackgo,wang2024boximator,zhang2024tora}, among which our work is particularly related to the tuning-free variants~\citep{qiu2024freetraj,ma2023trailblazer,DirectAVideo,jain2024peekaboo,Img2VidAnim-Zero}.
Yet, all of these methods focus on text-based generation. In contrast, our method uses an image-to-video model and thus can turn existing, real-world images into controllable videos.
Moreover, we can control camera motion by specifying trajectories of background regions. %



